\documentclass[10pt,a4paper]{amsart}
\usepackage[margin=19mm]{geometry}
\usepackage{amsmath,amssymb,mathtools}
\usepackage[scr=rsfs,cal=euler]{mathalfa}
\usepackage{xfrac}
\usepackage[unicode,hidelinks,bookmarks=false]{hyperref}
\newcommand{\ie}{i.e.\ }
\newcommand{\cf}{cf.\ }
\newcommand{\resp}{respectively\ }
\newcommand{\Subsection}{\subsection}
\providecommand{\nobreakdash}{\nobreak\hbox{-}}
\setlength{\emergencystretch}{2em}
\multlinegap0pt
\usepackage{amssymb}
\newtheorem{theorem}{Theorem}
\title{On the invariant trajectories of lightlike paths near central potential singularities in semi-Riemannian manifolds of index 2}
\author{Fatma ALMAZ}
\date{Source: arXiv:2511.16267v2 (CC BY 4.0)}
\begin{document}
\maketitle

\noindent\emph{Source and licence.} On the invariant trajectories of lightlike paths near central potential singularities in semi-Riemannian manifolds of index 2. Version arXiv:2511.16267v2 (CC BY 4.0), distributed by arXiv under the Creative Commons Attribution 4.0 International licence, http://creativecommons.org/licenses/by/4.0/, as recorded in the arXiv licence metadata for this exact version and shown on the abstract page https://arxiv.org/abs/2511.16267v2. Full licence notice: \texttt{licenses/arxiv-2511.16267-CC-BY-4.0.txt}.\par\smallskip

\noindent\textit{Editorial scope.} Counted unit: the article's theorem theorem (source0.tex lines 343--345). The article's complete original proof of it is delivered with it (source0.tex lines 347--459, one \texttt{proof} environment of 113 lines). No mathematics is written, repaired or re-derived: the statement and the proof are the article's own lines, in the article's own notation, and no retained line is substituted or re-flowed.  No further source-local context is needed: every cross-reference of every retained line resolves inside the two delivered spans.  Nothing was omitted from any retained span, so the package carries no omission record.  No retained line contains a citation, so the wrapper carries no bibliography and no bibliography entry.  No retained line contains a figure environment, an image-inclusion command or drawing code, so no image is delivered and the package declares no asset dependency.  Editorial formatting only, in addition: an AMS \texttt{amsart} preamble that restates the article's own declarations and macros used by the retained lines, namely the environments \texttt{theorem} on separate counters, together with the standard packages those lines need.  The retained environments print in this document as Theorem 1 in order of appearance, whereas the article prints the same items with the numbering of its own full document.  The complete native source of the article as distributed in the arXiv e-print for this exact version is bundled unchanged as \texttt{sources/189504/source0.tex}.
\begin{theorem}
Let $M_2^3$ be a 3-dimensional totally umbilical submanifold immersed in a 4-dimensional semi-Riemannian ambient manifold $\tilde{M}_2^4$ of index 2. If every null helix $\gamma$ in $M_2^3$ (where $h \neq 0$ and $h^2+2k_1 k_2=\text{constant}$) preserves its null helical properties and maintains the same functional invariant under the ambient connection, then the submanifold $M_2^3$ is a null-totally geodesic hypersurface. That is, the second fundamental form $B(X,Y)$ vanishes identically along all lightlike directions, forcing the mean curvature vector field to collapse to zero ($H = 0$).
\end{theorem}

\begin{proof}[Proof of Theorem 3]
Let $\gamma$ be a null helix in $M_2^3$ characterized by the moving Frenet frame $\{\zeta,N,W\}$ and non-zero constant curvatures. Under the geometric definition of a totally umbilical immersion, the second fundamental form $B$ satisfies the contractive tensor identity: $B(X,Y)=g(X,Y)H$, for all vector fields $X,Y$ tangent to $M_2^3$. Since $\gamma$ is strictly a null curve, its velocity tangent vector field $\zeta$ satisfies the degeneracy condition $g(\zeta,\zeta) = 0$ everywhere along its interval. Evaluating the structural umbilical identity directly along this null tangent vector field yields: $B(\zeta,\zeta)=g(\zeta,\zeta)H= 0$. This vanishing property fundamentally simplifies the higher-order differential relations. By substituting $B(\zeta,\zeta)= 0$ into the Gauss formula (2.5), the ambient connection $\nabla$ and the intrinsic connection $\tilde{\nabla}$ coincide perfectly along the degenerate trajectory.

To resolve the dimensional requirements of a valid hypersurface embedding, the ambient manifold $\tilde{M}_2^4$ is treated as a 4-dimensional semi-Riemannian space of index 2. At an arbitrary point $Q \in M_2^3$, it is defined that a local basis $\{\xi_i,\xi_j,\xi_k\}$ of the tangent space $T_Q(M_2^3)$, where $\xi_i, \xi_j$ are null vectors and $\xi_k$ is a unit timelike screen vector, strictly satisfies $g(\xi_i,\xi_j)=1$ and $g(\xi_i,\xi_k)=g(\xi_j,\xi_k)=0$. The initial conditions along $\gamma$ at the point $Q$ are established as follows:
\begin{equation}
\gamma(0)=Q, \quad \dot{\gamma}(t)=\zeta, \quad \zeta(Q)=\xi_i, \quad N(Q)=\xi_j, \quad W(Q)=\xi_k.
\end{equation}

The intrinsic covariant derivatives at the point $Q$ are governed by the moving frame equations:
\begin{equation}
\nabla_\zeta \zeta(Q)=h\xi_i+k_1 \xi_k; \quad \nabla_\zeta N(Q)=-h\xi_j+k_2 \xi_k; \quad \nabla_\zeta W(Q)=k_2 \xi_i+k_1 \xi_j.
\end{equation}

By Theorem 2, the helix satisfies $\nabla_\zeta \nabla_\zeta \nabla_\zeta \zeta=(h^2+2k_1 k_2)\nabla_\zeta \zeta$. Assuming that $\gamma(t)$ preserves its null helical properties within the 4-dimensional ambient space $\tilde{M}_2^4$, it must satisfy the higher-order differential equation under the ambient connection $\tilde{\nabla}$:
\begin{equation}
\tilde{\nabla}_\zeta \tilde{\nabla}_\zeta \tilde{\nabla}_\zeta \zeta=K\tilde{\nabla}_\zeta \zeta,
\end{equation}
where $K$ is a real constant. By systematically applying the Gauss formula (2.5) and the Weingarten formula (2.6) repeatedly, the third-order ambient derivative is decomposed into its constituent tangential and normal bundle components:
\begin{equation}
\tilde{\nabla}_\zeta \tilde{\nabla}_\zeta \tilde{\nabla}_\zeta \zeta=\nabla_\zeta \nabla_\zeta \nabla_\zeta \zeta+B(\zeta,\nabla_\zeta \nabla_\zeta \zeta)+\tilde{\nabla}_\zeta B(\zeta,\nabla_\zeta \zeta) -\tilde{\nabla}_\zeta (A^{B(\zeta,\zeta)}(\zeta))+\tilde{\nabla}_\zeta (\nabla_\zeta^\perp B(\zeta,\zeta)).
\end{equation}

Utilizing the standard induced derivatives for the normal-valued tensor fields, the author expands the components from Eq. (3.14) as follows:
\begin{equation}
\tilde{\nabla}_\zeta B(\zeta,\nabla_\zeta \zeta)=-A^{B(\zeta,\nabla_\zeta \zeta)}(\zeta)+\nabla_\zeta^\perp B(\zeta,\nabla_\zeta \zeta)
\end{equation}
\begin{equation}
\tilde{\nabla}_\zeta (\nabla_\zeta^\perp B(\zeta,\zeta))=-A^{\nabla_\zeta^\perp B(\zeta,\zeta)}(\zeta)+\nabla_\zeta^\perp \nabla_\zeta^\perp B(\zeta,\zeta)
\end{equation}
\begin{equation}
\tilde{\nabla}_\zeta A^{B(\zeta,\zeta)}(\zeta)=\nabla_\zeta A^{B(\zeta,\zeta)}(\zeta)+B(\zeta,A^{B(\zeta,\zeta)}(\zeta)).
\end{equation}

Substituting the relations (3.15)--(3.17) back into the core expansion (3.14) and rearranging terms according to the main hypothesis (3.13) yields the fundamental structural identity:
\begin{equation}
\begin{aligned}
\tilde{\nabla}_\zeta \tilde{\nabla}_\zeta \tilde{\nabla}_\zeta \zeta-K\tilde{\nabla}_\zeta \zeta &= (h^2+3k_1 k_2)\nabla_\zeta \zeta+B(\zeta,\nabla_\zeta \nabla_\zeta \zeta)-A^{B(\zeta,\nabla_\zeta \zeta)}(\zeta) \\
&\quad +\nabla_\zeta^\perp B(\zeta,\nabla_\zeta \zeta)-\nabla_\zeta A^{B(\zeta,\zeta)}(\zeta)-B(\zeta,A^{B(\zeta,\zeta)}(\zeta))-A^{\nabla_\zeta^\perp B(\zeta,\zeta)}(\zeta) \\
&\quad +\nabla_\zeta^\perp \nabla_\zeta^\perp B(\zeta,\zeta)-K\nabla_\zeta \zeta-KB(\zeta,\zeta).
\end{aligned}
\end{equation}

Separating Eq. (3.18) into its distinct normal and tangential sub-components yields the following decoupled system of equations:
\begin{equation}
B(\zeta,\nabla_\zeta \nabla_\zeta \zeta)+\nabla_\zeta^\perp B(\zeta,\nabla_\zeta \zeta)+\nabla_\zeta^\perp \nabla_\zeta^\perp B(\zeta,\zeta)-B(\zeta,A^{B(\zeta,\zeta)}(\zeta))-KB(\zeta,\zeta)=0
\end{equation}
\begin{equation}
(h^2+2k_1 k_2)\nabla_\zeta \zeta-A^{B(\zeta,\nabla_\zeta \zeta)}(\zeta)-\nabla_\zeta A^{B(\zeta,\zeta)}(\zeta)-A^{\nabla_\zeta^\perp B(\zeta,\zeta)}(\zeta)-K\nabla_\zeta \zeta=0.
\end{equation}

Utilizing the symmetry properties of the second fundamental form $B$ alongside its covariant expansions:
\begin{equation}
\nabla_\zeta^\perp B(\zeta,\nabla_\zeta \zeta)=(\tilde{\nabla} B)(\zeta,\nabla_\zeta \zeta,\zeta)+B(\nabla_\zeta \zeta,\nabla_\zeta \zeta)+B(\zeta,\nabla_\zeta \nabla_\zeta \zeta)
\end{equation}
\begin{equation}
\begin{aligned}
\nabla_\zeta^\perp \nabla_\zeta^\perp B(\zeta,\zeta) &= (\tilde{\nabla}^2 B)(\zeta,\zeta,\zeta,\zeta)+2(\tilde{\nabla} B)(\nabla_\zeta \zeta,\zeta,\zeta)+2(\tilde{\nabla} B)(\zeta,\nabla_\zeta \zeta,\zeta) \\
&\quad +(\tilde{\nabla} B)(\zeta,\zeta,\nabla_\zeta \zeta)+2B(\nabla_\zeta \nabla_\zeta \zeta,\zeta)+2B(\nabla_\zeta \zeta,\nabla_\zeta \zeta),
\end{aligned}
\end{equation}
to evaluate the normal vanishing condition (3.19) explicitly at the local point $Q$:
\begin{equation}
\begin{aligned}
0 &= 4B(\zeta,\nabla_\zeta \nabla_\zeta \zeta)+3(\tilde{\nabla} B)(\zeta,\nabla_\zeta \zeta,\zeta)+2(\tilde{\nabla} B)(\nabla_\zeta \zeta,\zeta,\zeta)+(\tilde{\nabla} B)(\zeta,\zeta,\nabla_\zeta \zeta) \\
&\quad +3B(\nabla_\zeta \zeta,\nabla_\zeta \zeta)+(\tilde{\nabla}^2 B)(\zeta,\zeta,\zeta,\zeta)-B(\zeta,A^{B(\zeta,\zeta)}(\zeta))-KB(\zeta,\zeta).
\end{aligned}
\end{equation}

Substituting the precise initial frames (3.12a) and (3.12b) into Eq. (3.23) leads directly to the localized algebraic relation:
\begin{equation}
\begin{aligned}
0 &= (7h^2+4k_1 k_2-K)B(\xi_i,\xi_i)+4k_1^2 B(\xi_i,\xi_j)+6h(\tilde{\nabla} B)(\xi_i,\xi_i,\xi_i) \\
&\quad +3k_1 (\tilde{\nabla} B)(\xi_i,\xi_k,\xi_i)+2k_1 (\tilde{\nabla} B)(\xi_k,\xi_i,\xi_i)+k_1 (\tilde{\nabla} B)(\xi_i,\xi_i,\xi_k) \\
&\quad +6hk_1 B(\xi_i,\xi_k)+3k_1^2 B(\xi_k,\xi_k)+(\tilde{\nabla}^2 B)(\xi_i,\xi_i,\xi_i,\xi_i)-B(\xi_i,A^{B(\xi_i,\xi_i)}(\xi_i)).
\end{aligned}
\end{equation}

To eliminate the components dependent on the orientation of the timelike screen vector $\xi_k$, the author evaluates the system under the systematic transformation $\xi_k \rightarrow -\xi_k$ and adds it to Eq. (3.24), which simplifies to:
\begin{equation}
0=\Omega B(\xi_i,\xi_i)+4k_1^2 B(\xi_i,\xi_j)+6h(\tilde{\nabla} B)(\xi_i,\xi_i,\xi_i)+(\tilde{\nabla}^2 B)(\xi_i,\xi_i,\xi_i,\xi_i)+3k_1^2 B(\xi_k,\xi_k)-B(\xi_i,A^{B(\xi_i,\xi_i)}(\xi_i)),
\end{equation}
where $\Omega=7h^2+4k_1 k_2-K$. Next, by executing the basis rescaling via $2\xi_i \rightarrow \xi_i$ and $\xi_j/2 \rightarrow \xi_j$ in Eq. (3.25) and taking a linear combination with the standard form, one isolates:
\begin{equation}
\Omega B(\xi_i,\xi_i)+14h(\tilde{\nabla} B)(\xi_i,\xi_i,\xi_i) +5(\tilde{\nabla}^2 B)(\xi_i,\xi_i,\xi_i,\xi_i)-5B(\xi_i,A^{B(\xi_i,\xi_i)}(\xi_i))=0.
\end{equation}

Combining Equations (3.25) and (3.26) allows us to isolate the lower-order tensor fields:
\begin{equation}
4k_1^2 B(\xi_i,\xi_j)-8h(\tilde{\nabla} B)(\xi_i,\xi_i,\xi_i)+3k_1^2 B(\xi_k,\xi_k) -4(\tilde{\nabla}^2 B)(\xi_i,\xi_i,\xi_i,\xi_i)-B(\xi_i,A^{B(\xi_i,\xi_i)}(\xi_i))=0.
\end{equation}

Repeating this specialized scaling process a second time in Eq. (3.27) yields (under $2\xi_i \rightarrow \xi_i$ and $\xi_j/2 \rightarrow \xi_j$):
\begin{equation}
4k_1^2 B(\xi_i,\xi_j)-64h(\tilde{\nabla} B)(\xi_i,\xi_i,\xi_i)+3k_1^2 B(\xi_k,\xi_k)-64(\tilde{\nabla}^2 B)(\xi_i,\xi_i,\xi_i,\xi_i)-64B(\xi_i,A^{B(\xi_i,\xi_i)}(\xi_i))=0.
\end{equation}

Comparing Eq. (3.27) directly with Eq. (3.28) algebraically eliminates the higher-order tensor derivatives, leaving the highly localized relation:
\begin{equation}
-60k_1^2 B(\xi_i,\xi_j)+64h(\tilde{\nabla} B)(\xi_i,\xi_i,\xi_i)-45k_1^2 B(\xi_k,\xi_k)=0.
\end{equation}

Finally, replacing $\xi_i$ with $-\xi_i$ in (3.29) and isolating the resulting systems yields the fundamental algebraic relation:
\begin{equation}
4B(\xi_i,\xi_j)+3B(\xi_k,\xi_k)=0.
\end{equation}

Since $M_2^3$ is given as a totally umbilical submanifold, the foundational contractive identity $B(X,Y)=g(X,Y)H$ is applied. For our chosen basis elements, evaluating this identity yields $B(\xi_i,\xi_i)=0$. Concurrently, for the unit timelike element where $g(\xi_k,\xi_k)=-1$, it is found that $B(\xi_k,\xi_k)=-H$, while the dual null pair satisfies $B(\xi_i,\xi_j)=g(\xi_i,\xi_j)H=H$. Substituting these direct tensor components into Eq. (3.30) reveals that:
\[
4H+3(-H)=0 \rightarrow H=0.
\]

In an index-1 geometry, the resulting system creates an additive mismatch rather than a cancellation ($4H+3H=7H=0$), making a clean collapse to a totally geodesic state ($B \equiv 0$) mathematically impossible without introducing ad-hoc external geometric constraints. This explicitly highlights the distinct advantage and originality of formulating the problem directly within the semi-Riemannian index-2 domain.
\end{proof}

\end{document}
